\documentclass[a4paper,10pt,twocolumn]{myjarticle2}

\usepackage[dvipdfmx]{graphicx}
\usepackage{url}

\input{header}

\renewcommand{\figurename}{Figure}
\renewcommand{\tablename}{Table}
\renewcommand{\refname}{References}

\pagestyle{empty}

\setlength{\topmargin}{6mm}
\setlength{\oddsidemargin}{-4mm}
\setlength{\evensidemargin}{-4mm}
\setlength{\textwidth}{170mm}
\setlength{\headsep}{0pt}
\setlength{\headheight}{0pt}
\setlength{\textheight}{235mm}
\setlength{\columnsep}{5mm}

\newcommand{\Hraw}{H_{\mathrm{raw}}}
\newcommand{\Hstar}{H^{\ast}}

\begin{document}
\twocolumn[%
\vspace{-20mm}
\begin{center}
{\LARGE\bf Membership Specification and Sensitivity Reporting for
Fuzzy-Entropy-Based LLM Emotion Scoring}
\vspace{2mm}\par\noindent
{\large
$\bigcirc$ Naruki Shirahama$^{1}$, Yuma Yoshimoto$^{2}$, Naofumi Nakaya$^{3}$,\\
Junji Nishino$^{4}$, and Satoshi Watanabe$^{5}$}
\vspace{2mm}\par\noindent
{\small
$^{1}$Shimonoseki City University, Shimonoseki, Yamaguchi, Japan\\
$^{2}$NIT (KOSEN), Kitakyushu College, Kitakyushu, Fukuoka, Japan\\
$^{3}$Juntendo University, Urayasu, Chiba, Japan\\
$^{4}$Tamagawa University, Machida, Tokyo, Japan\\
$^{5}$Shizuoka Institute of Science and Technology, Fukuroi, Shizuoka, Japan\\
Contact: \url{nshirahama@ieee.org}}
\vspace{2mm}\par\noindent
\begin{minipage}{150mm}
{\small
{\bf Abstract:}
Fuzzy entropy has been used as an ambiguity descriptor for continuous emotion
scores produced by large language models. However, the resulting entropy value
is not determined by the score alone; it also depends on the
membership-function specification used to map a 0--100 score into Low, Medium,
and High fuzzy memberships. This short note proposes a membership specification
card and a minimal sensitivity-reporting protocol for fuzzy-entropy-based LLM
emotion scoring. The proposed report records the score range, labels,
breakpoints, membership family, smoothness, endpoint behavior, entropy formula,
family-specific normalization, and version identifier. Using three score-axis
membership families--piecewise linear, Zadeh-type S/Z/pi, and
sigmoid-composed S/Z/pi--we illustrate how normalized entropy-response curves
$\Hstar(x)$ and family-pair sensitivity summaries can be reported without
introducing model rankings or persona-temperature analyses. The aim is not to
select a universally correct membership function, but to make fuzzy-entropy
reports reproducible and auditable.
\par\noindent
}
\end{minipage} \par\noindent
\end{center}
]

\section{Introduction}

Continuous emotion scores produced by large language models (LLMs) are often
reported on a bounded scale, for example from 0 to 100. Fuzzy-set methods offer
a convenient way to describe such scores using linguistic memberships such as
Low, Medium, and High. A scalar fuzzy entropy can then be used as an ambiguity
descriptor: scores near a single dominant label have low entropy, while scores
near transitions among labels have higher entropy.
Prior EQ-Fuzzy work used fuzzy entropy as one LLM emotion-scoring descriptor
[4]; this note uses that context only as motivation for reproducible reporting.

This descriptor is useful only when its membership specification is visible.
A score value does not determine its fuzzy entropy by itself. The entropy also
depends on the score range, breakpoints, membership-function family,
smoothness, endpoint behavior, and normalization convention. Two studies can
therefore use the same LLM score and the same entropy formula but report
different entropy values because their membership functions differ.

This paper proposes a compact reporting unit for fuzzy-entropy-based LLM
emotion scoring. The contribution is methodological rather than evaluative:
it does not evaluate which LLM is best, analyze persona-temperature effects,
provide human-grounded multilingual validation, or position EQ-Fuzzy against
existing benchmarks. Instead, it specifies what should be reported when fuzzy
entropy is used as a score-axis descriptor.

\section{Fuzzy Entropy and Membership Dependence}

Let $x \in [0,100]$ be a continuous emotion score. A membership specification
maps $x$ to three membership degrees
$\mu_L(x), \mu_M(x), \mu_H(x) \in [0,1]$ for Low, Medium, and High. The
artifact used in this note computes a Shannon-type membership-overlap entropy
\begin{equation}
\Hraw(x)=-\sum_{\ell \in \{L,M,H\}}\mu_\ell(x)\log_2\mu_\ell(x),
\end{equation}
where $0\log_2 0:=0$. For comparison across membership families, we report a
family-normalized response
\begin{equation}
\Hstar(x) = \Hraw(x) / H_{\max},
\end{equation}
where $H_{\max}$ is estimated for that family on the documented score grid.

The dependence on membership functions is not incidental. Zadeh's fuzzy sets
are defined through membership grades over a universe of discourse [1]. The
De Luca--Termini formulation is part of the fuzzy-entropy literature [2], but
the concrete descriptor reported here is the membership-overlap entropy in
Eq. (1). Recent FSS discussion of membership functions also emphasizes that
values of the universe and the shape of the membership function jointly affect
interpretation [3]. For LLM emotion scoring, this means that an entropy number
should be read as ``entropy under a documented membership specification.''

\section{Membership Specification Card}

Table 1 shows family-specific fields in the proposed membership specification
card. The card records the minimum information needed to reproduce and audit a
fuzzy-entropy descriptor over the score axis. The exact version identifiers and
configuration hashes are retained in the machine-readable artifact; the
manuscript uses short family aliases to keep the report readable in proceedings
format.

\begin{table}
\centering
\caption{Family-specific fields in the membership specification card.}
\label{tab:membership-card}
\footnotesize
\setlength{\tabcolsep}{2pt}
\begin{tabular}{|p{18mm}|p{16mm}|p{20mm}|p{10mm}|}
\hline
Family & Smoothness & Endpoint handling & $H_{\max}$ \\
\hline
linear & $C^0$ piecewise & exact endpoints & 1.0577 \\
\hline
Zadeh S/Z/pi & $C^1$ quadratic & exact endpoints & 1.0524 \\
\hline
sigmoid S/Z/pi & smooth transition & clipped interior values & 1.0184 \\
\hline
\end{tabular}
\end{table}

The three families illustrate common membership-function design choices.
Piecewise-linear memberships are interpretable but nonsmooth at transition
points. Zadeh-type S/Z/pi families smooth local transitions. Sigmoid-composed
variants document score-axis warping before S/Z/pi mapping. These choices
determine where a 0--100 score is treated as crisp or transition-like.

The complete card should also state the score range $[0,100]$, score-grid
resolution, labels, Low/Medium and Medium/High breakpoints, entropy formula,
normalization rule, and version identifier. These fields are not decorative.
For example, changing a breakpoint changes which scores are treated as
transition scores; changing endpoint behavior changes whether the extreme
scores have exactly zero ambiguity; changing $H_{\max}$ changes comparability
across families.

\section{Entropy-Response and Sensitivity Reporting}

Figure 1 illustrates the reporting protocol. The upper panels
show Low, Medium, and High memberships over the score axis for three membership
families. The lower panel shows the corresponding normalized
entropy-response curves $\Hstar(x)$. The figure is deliberately score-axis-only:
it contains no model, provider, persona, temperature, benchmark, language-drift,
or target-shift results.

\begin{figure}
\centering
\includegraphics[width=\linewidth]{figures/figure1_membership_and_entropy_response.png}
\caption{Score-axis-only membership functions and normalized entropy-response
curves. The figure reports membership-family behavior over the 0--100 score
axis only and contains no model, provider, persona, temperature, benchmark, or
target-shift results.}
\label{fig:response}
\end{figure}

A minimal sensitivity report compares unordered family pairs over the same
score grid. Table 2 shows that the largest absolute differences in
$\Hstar(x)$ occur in the Low--Medium transition region. These values are not
model effects. They are score-axis sensitivity summaries under alternative
membership specifications.

\begin{table}
\centering
\caption{Score-axis sensitivity summary over the documented 0--100 grid.}
\label{tab:sensitivity}
\footnotesize
\setlength{\tabcolsep}{2pt}
\begin{tabular}{|p{23mm}|p{11mm}|p{12mm}|p{18mm}|}
\hline
Pair & Max $|\Delta\Hstar|$ & Mean $|\Delta\Hstar|$ & max region \\
\hline
linear--Zadeh S/Z/pi & 0.343 & 0.110 & Low--Medium \\
\hline
linear--sigmoid-composed & 0.481 & 0.158 & Low--Medium \\
\hline
Zadeh S/Z/pi--sigmoid-composed & 0.167 & 0.048 & Low--Medium \\
\hline
\end{tabular}
\end{table}

The minimum reporting checklist is: (i) score range and grid resolution;
(ii) labels and breakpoints; (iii) family ID, smoothness, and endpoint
convention; (iv) entropy equation and family-specific $H_{\max}$; and
(v) score-axis sensitivity against at least one plausible alternative family.
This is intentionally smaller than a full experiment table. It is meant to
make the descriptor auditable before any downstream LLM comparison is
interpreted.

\section{Discussion and Conclusion}

The proposed card does not identify a universally correct membership function.
In many LLM emotion-scoring settings, several membership families may be
defensible. The practical requirement is that the chosen family be specified
and that nearby plausible choices be sensitivity-checked. Without that record,
fuzzy entropy values can appear more self-contained than they really are.

This short note reframes fuzzy entropy for LLM emotion scoring as a
specification-dependent descriptor. A reproducible report should state the
membership-function family, breakpoints, smoothness, endpoint handling,
entropy formula, and family-specific normalization, then show the
score-axis entropy response and minimal sensitivity summary. This small
reporting unit is intended to support later experiments without importing
model rankings, persona-temperature decompositions, human-alignment rankings,
or benchmark-positioning claims into the membership specification itself.

\vspace{1mm}
\noindent{\large\bf Acknowledgment}\par\noindent

This work is part of the EQ-Fuzzy research line. This work was partially
supported by JSPS Grants-in-Aid for Scientific Research (JSPS KAKENHI), Grant
Numbers 26K14983 and 24K13350. The present manuscript uses only generated
score-axis membership-specification artifacts and does not use API calls or raw
JSONL outputs.

\vspace{1mm}
\noindent{\large\bf References}\par\noindent
[1] L. A. Zadeh, ``Fuzzy sets,'' \emph{Information and Control}, vol. 8,
no. 3, pp. 338--353, 1965.
\par\noindent
[2] A. De Luca and S. Termini, ``A definition of a nonprobabilistic entropy in
the setting of fuzzy sets theory,'' \emph{Information and Control}, vol. 20,
no. 4, pp. 301--312, 1972.
\par\noindent
[3] M. Umano, ``Some considerations on membership functions of fuzzy sets,''
\emph{Proceedings of the Fuzzy System Symposium}, vol. 41, no. 0, pp. 70--74,
2025, \url{https://doi.org/10.14864/fss.41.0_70}.
\par\noindent
[4] N. Shirahama, Y. Yoshimoto, N. Nakaya, and S. Watanabe, ``Mathematical
framework for characterizing emotional individuality in large language models:
Temperature control, fuzzy entropy, and persona-based diversity analysis,''
\emph{Mathematics}, vol. 14, no. 7, Art. no. 1224, 2026,
doi:10.3390/math14071224.

%%%%% Contact Information%%%%%
\vspace{1mm}
\noindent{\large\bf Contact Information}\par\noindent
\begin{tabular}{l}
Naruki Shirahama \\
Shimonoseki City University\\
E-mail: {\tt  nshirahama@ieee.org}
\end{tabular}

\end{document}
